%% 1-10 時こくと 時間（第11回で 時計の よみかたの 段を ふやした）
%% 第1巻のきまり：15 pt・行送り28 pt、分かち書き（文節の間に半角の空白）、漢字はすべてルビ、式は全角の＋－＝で書く
%% この章の部品（名前は \bkqi…）
\makeatletter
% 時計のまわりに 5、10、15 … と書く  \bkqifives{(x,y)}{時計の半径}{いくつまで（5の倍数）}
\newcommand\bkqifives[3]{\begin{scope}[shift={#1}]
  \foreach \k in {1,...,12} {\pgfmathtruncatemacro\bkqi@v{5*\k}%
    \ifnum\bkqi@v>#3\relax\else
      \node[font=\sffamily\bfseries\fontsize{8.5}{9}\selectfont, text=bkBlue!85!black, inner sep=0pt] at ({90-30*\k}:#2+4.6) {\bkqi@v};\fi}
\end{scope}}
% ストップウォッチ  \bkqiwatch{(x,y)}{半径}{秒}
\newcommand\bkqiwatch[3]{\begin{scope}[shift={#1}]
  \fill[bkGray!70] (-1.4,#2) rectangle (1.4,#2+3.4);
  \fill[bkGray!85, rounded corners=0.8mm] (-3.4,#2+3) rectangle (3.4,#2+5.4);
  \begin{scope}[rotate=-42]\fill[bkGray!70, rounded corners=0.5mm] (-1.3,#2) rectangle (1.3,#2+2.8);\end{scope}
  \filldraw[fill=bkBlue!25, draw=bkBlue!70!black, line width=1.2pt] (0,0) circle (#2+1.6);
  \filldraw[fill=white, draw=bkInk!50, line width=0.5pt] (0,0) circle (#2);
  \foreach \m in {0,...,59} {\pgfmathsetmacro\bkqi@a{90-6*\m}\pgfmathtruncatemacro\bkqi@big{mod(\m,5)==0}
    \ifnum\bkqi@big=1 \draw[bkInk, line width=0.8pt] (\bkqi@a:#2*0.94) -- (\bkqi@a:#2*0.8);
    \else\draw[bkInk!70, line width=0.35pt] (\bkqi@a:#2*0.94) -- (\bkqi@a:#2*0.87);\fi}
  \pgfmathsetmacro\bkqi@fs{#2*0.42}%
  \foreach \k in {1,...,12} {\pgfmathtruncatemacro\bkqi@v{5*\k}%
    \node[font=\sffamily\bfseries\fontsize{\bkqi@fs}{\bkqi@fs}\selectfont, text=bkInk, inner sep=0pt] at ({90-30*\k}:#2*0.62) {\bkqi@v};}
  \draw[bkRed, line width=1.1pt, line cap=round] (0,0) -- ({90-6*#3}:#2*0.86);
  \fill[bkRed] (0,0) circle (0.9);
\end{scope}}
% いえ  \bkqihouse{(x,y)}{はば}{かべの色}   (x,y) は下のまん中
\newcommand\bkqihouse[3]{\begin{scope}[shift={#1}]
  \fill[#3] (-#2/2,0) rectangle (#2/2,#2*0.62);
  \fill[bkRed!75] (-#2/2-1.2,#2*0.6) -- (0,#2*1.08) -- (#2/2+1.2,#2*0.6) -- cycle;
  \fill[bkStem!80] (-#2*0.12,0) rectangle (#2*0.12,#2*0.36);
  \fill[white] (#2*0.2,#2*0.26) rectangle (#2*0.38,#2*0.44);
\end{scope}}
% 「□時□分」の書きこみ（答えだけ。行の右によせる）  \bkqijifun
\newcommand\bkqiright[1]{\unskip\nobreak\hfil\penalty50\hskip1em\null\nobreak\hfill\mbox{#1}\par}
\newcommand\bkqijifun{\bkqiright{{\sffamily\bfseries\small\color{bkOrange}こたえ}\enspace
  \bkkakikomi\,\R{時}{じ}\enspace\bkkakikomi\,\R{分}{ふん}}}
% 「□分□秒」などの書きこみ  \bkqitwo{単位1}{単位2}
\newcommand\bkqitwo[2]{\bkqiright{{\sffamily\bfseries\small\color{bkOrange}こたえ}\enspace
  \bkkakikomi\,#1\enspace\bkkakikomi\,#2}}
\makeatother

\bkchapter{1}{10}{\R{時}{じ}こくと \R{時間}{じ|かん}}{\bklevel{1}}{\R{時}{じ}こくと \R{時間}{じ|かん}}
\bkmichishirube{1-2,1-5,1-6}{1-15,1-18,2-10}

\begin{goal}
\begin{itemize}[label={\color{bkTheme}●}, leftmargin=1.4\zw, itemsep=0pt, topsep=0pt]
\item \mbox{\R{時計}{と|けい}}を よんで、 \R{時}{じ}こくが いえる。
\item 「\R{時}{じ}こく」と 「\mbox{\R{時間}{じ|かん}}」の ちがいが わかる。
\item 60で くり\R{上}{あ}がる（1-7の くり\R{上}{あ}がりと おなじ） \mbox{\R{時間}{じ|かん}}の けいさんが できる。
\end{itemize}
\end{goal}

\begin{hatena}
\begin{bkfigure}
\begin{tikzpicture}[x=1mm, y=1mm]
  % A：おやつの 3時
  \fill[white, rounded corners=3mm] (0,0) rectangle (59,54);
  \node[circle, fill=bkYellow, text=white, font=\sffamily\bfseries, minimum size=6.5mm, inner sep=0pt] at (6,48) {A};
  \node[anchor=west, font=\small, text=bkInk] at (11,48) {おやつは \mbox{3\R{時}{じ}}};
  \bkclock{(21,21)}{15}{3}{0}
  \bkplate{(48,10)}{18}
  \foreach \x/\y in {44/12,51/13,47.5/16.5} {\bkcookie{(\x,\y)}{3}}
  % B：おじいちゃんの いえまで 3時間
  \begin{scope}[shift={(63,0)}]
  \fill[white, rounded corners=3mm] (0,0) rectangle (59,54);
  \node[circle, fill=bkYellow, text=white, font=\sffamily\bfseries, minimum size=6.5mm, inner sep=0pt] at (6,48) {B};
  \node[anchor=west, font=\small, text=bkInk] at (11,48) {くるまで \mbox{3\R{時間}{じ|かん}}};
  \fill[bkGray!35] (3,14) rectangle (56,18);
  \draw[white, line width=0.8pt, dash pattern=on 2.5pt off 2pt] (3,16) -- (56,16);
  \bkqihouse{(10,18)}{11}{bkBlue!35}
  \bkqihouse{(48,18)}{11}{bkGreen!35}
  \bkcar{(24,16.2)}{0.42}{bkRoute6}
  \node[anchor=north, font=\sffamily\bfseries\fontsize{9}{11}\selectfont, text=bkInk] at (10,12.5) {いえ};
  \node[anchor=north, font=\sffamily\bfseries\fontsize{9}{11}\selectfont, text=bkInk, align=center] at (46,12.5) {おじいちゃん\\の いえ};
  \bkbrace{(10,31)}{(48,31)}{\mbox{3\R{時間}{じ|かん}}}{1}
  \end{scope}
\end{tikzpicture}
\end{bkfigure}
\noindent A の 「\mbox{3\R{時}{じ}}」と、 B の 「\mbox{3\R{時間}{じ|かん}}」。\\
どちらにも 「3」と 「\R{時}{じ}」が ついて いるね。\\
でも、 おなじ ことかな？ どこが ちがうのかな？
\end{hatena}

\clearpage
%% ---------------- p.99 やってみよう ----------------
\begin{yattemiyou}[\mbox{\R{時計}{と|けい}}の はりを うごかそう]
\begin{enumerate}[label={\bkcircled[bkGreen]{\arabic*}}, leftmargin=2.2\zw, itemsep=2pt, topsep=0pt]
\item \mbox{\R{時計}{と|けい}}の はりを \mbox{3\R{時}{じ}}に あわせよう。
\item \R{長}{なが}い はりを ぐるっと 1まわり させよう。 みじかい はりは どこへ いったかな？
\item その あいだに \R{長}{なが}い はりは、 \R{小}{ちい}さい \R{目}{め}もりを いくつ すすんだかな？ まわりの 5、 10、 15… を つかって かぞえよう。
\end{enumerate}
\begin{bkfigure}
\begin{tikzpicture}[x=1mm, y=1mm]
  \bkclock{(17,4)}{16}{3}{0}
  \draw[bkInk, line width=2.6pt, line cap=round] (2,-18) -- (9,-18);
  \node[anchor=west, font=\sffamily\bfseries\footnotesize] at (11,-18) {みじかい はり};
  \draw[bkRed!85!black, line width=1.4pt, line cap=round] (0,-25) -- (9,-25);
  \node[anchor=west, font=\sffamily\bfseries\footnotesize, text=bkRed!85!black] at (11,-25) {\R{長}{なが}い はり};
  \draw[bk arrow, bkGreen] (39,4) -- (54,4);
  \node[font=\sffamily\bfseries\footnotesize, text=bkGreen!70!black] at (46.5,9) {1まわり};
  \bkclock{(90,0)}{19}{4}{0}
  \bkqifives{(90,0)}{19}{60}
  \draw[bkGreen, line width=1.5pt, -{Stealth[length=6pt]}] ($(90,0)+(84:29.5)$) arc (84:-258:29.5);
\end{tikzpicture}
\end{bkfigure}
\noindent\hspace*{1\zw}②\enspace みじかい はりは\enspace\bkkakikomi\enspace の ところ\hspace{2.5\zw}③\enspace\bkkakikomi\enspace こ
\end{yattemiyou}

\bkfukidashi{sora}{\R{長}{なが}い はりが 1まわり すると、 みじかい はりは つぎの \mbox{\R{数字}{すう|じ}}まで、 1つだけ すすむんだね！}

\clearpage
%% ---------------- 数学の道具（第11回：時計の よみかたを 1段ずつ ふやした） ----------------
\bksection{doguu}{\mbox{\R{数学}{すう|がく}}の \mbox{\R{道具}{どう|ぐ}}}

\bksubsection{1}{\mbox{\R{時計}{と|けい}}の しくみ}

\label{bk:1-10:shikumi}\mbox{\R{時計}{と|けい}}には、 はりが 2\R{本}{ほん}と、 \R{数字}{すう|じ}が 12こ、 \R{小}{ちい}さい \R{目}{め}もりが あります。

\begin{bkfigure}
\begin{tikzpicture}[x=0.78mm, y=0.78mm]
  \bkclock{(30,0)}{25}{10}{10}
  \draw[bkInk, line width=0.8pt] (5,15) -- (-14,22);
  \node[anchor=east, font=\sffamily\bfseries\footnotesize, align=right] at (-14,22) {みじかい はり\\（\R{時}{じ}を よむ）};
  \draw[bkRed!85!black, line width=0.8pt] (45,15) -- (64,22);
  \node[anchor=west, font=\sffamily\bfseries\footnotesize, text=bkRed!85!black, align=left] at (64,22) {\R{長}{なが}い はり\\（\R{分}{ふん}を よむ）};
  \draw[bkInk, line width=0.8pt] (52.5,-11) -- (64,-14);
  \node[anchor=west, font=\sffamily\bfseries\footnotesize, align=left] at (64,-14) {\R{小}{ちい}さい \R{目}{め}もり\\ぜんぶで 60こ};
  \draw[bkInk, line width=0.8pt] (12,-8) -- (-14,-14);
  \node[anchor=east, font=\sffamily\bfseries\footnotesize, align=right] at (-14,-14) {\R{数字}{すう|じ}は 1から 12};
\end{tikzpicture}
\end{bkfigure}
\begin{itemize}[label={\color{bkTheme}●}, leftmargin=1.4\zw, itemsep=0pt, topsep=2pt]
\item みじかい はりは ゆっくり うごき、 「なん\R{時}{じ}」を おしえる。
\item \R{長}{なが}い はりは はやく うごき、 「なん\R{分}{ぷん}」を おしえる。
\item \R{数字}{すう|じ}と \R{数字}{すう|じ}の あいだには、 \R{小}{ちい}さい \R{目}{め}もりの あいだが 5つ ある。
\end{itemize}

\noindent{\sffamily\bfseries\color{bkTheme}いろいろな \R{時計}{と|けい}}\quad かたちは ちがっても、 よみかたは おなじ。
\begin{bkfigure}
\begin{tikzpicture}[x=0.82mm, y=0.82mm]
  % めざまし時計
  \fill[bkRed!70] (-9,9) circle (3.2) (9,9) circle (3.2);
  \draw[bkInk!70, line width=1pt] (-6,-11) -- (-8,-14) (6,-11) -- (8,-14);
  \fill[bkRed!70] (0,0) circle (12.5);
  \bkclock{(0,0)}{11}{7}{0}
  \node[anchor=north, font=\sffamily\bfseries\footnotesize] at (0,-15) {めざまし\R{時計}{ど|けい}};
  % うで時計
  \begin{scope}[shift={(45,0)}]
  \fill[bkStem!70, rounded corners=1.5mm] (-5,-16) rectangle (5,16);
  \fill[bkGray!60] (0,0) circle (10.6);
  \bkclock{(0,0)}{9.5}{3}{40}
  \node[anchor=north, font=\sffamily\bfseries\footnotesize] at (0,-17) {うで\R{時計}{ど|けい}};
  \end{scope}
  % えきの 時計
  \begin{scope}[shift={(90,0)}]
  \fill[bkGray!70] (-1.2,12) rectangle (1.2,17);
  \fill[bkInk!80] (0,0) circle (13);
  \bkclock{(0,0)}{11.5}{10}{25}
  \node[anchor=north, font=\sffamily\bfseries\footnotesize] at (0,-15) {えきの \R{時計}{と|けい}};
  \end{scope}
\end{tikzpicture}
\end{bkfigure}
\noindent\R{左}{ひだり}から \mbox{7\R{時}{じ}}、 \mbox{3\R{時}{じ}40\R{分}{ぷん}}、 \mbox{10\R{時}{じ}25\R{分}{ふん}}。 よみかたは つぎから。



\clearpage
\bksubsection{2}{みじかい はりで 「なん\R{時}{じ}」}

\label{bk:1-10:yomu}みじかい はりは、 8から 9まで、 \mbox{1\R{時間}{じ|かん}} かけて ゆっくり すすみます。 その とちゅうを \R{見}{み}て みよう。

\begin{bkfigure}
\begin{tikzpicture}[x=1mm, y=1mm]
  \foreach \i/\m/\t in {0/0/{8\R{時}{じ}},1/15/{8\R{時}{じ}15\R{分}{ふん}},2/30/{8\R{時}{じ}30\R{分}{ぷん}},3/45/{8\R{時}{じ}45\R{分}{ふん}}} {
    \bkclock{(\i*26+11,0)}{11}{8}{\m}
    \node[anchor=north, font=\sffamily\bfseries\scriptsize] at (\i*26+11,-12.5) {\t};}
  \bkclock{(4*26+11,0)}{11}{9}{0}
  \node[anchor=north, font=\sffamily\bfseries\scriptsize, text=bkTheme] at (4*26+11,-12.5) {9\R{時}{じ}};
  \fill[bkBlue!14, rounded corners=1.5mm] (-1,-21.5) rectangle (99,-27.5);
  \node[font=\sffamily\bfseries\footnotesize, text=bkBlue!80!black] at (49,-24.5) {みじかい はりが 8と 9の あいだ → まだ 8\R{時}{じ}};
  \node[font=\sffamily\bfseries\footnotesize, text=bkTheme] at (115,-24.5) {9に ついた};
\end{tikzpicture}
\end{bkfigure}
\noindent みじかい はりが 8と 9の あいだに ある あいだは、 まだ 9\R{時}{じ}に なって いないので 「8\R{時}{じ}なん\R{分}{ぷん}」です。 9に ついた とき、 はじめて 9\R{時}{じ}に なります。

\begin{taisetsu}
みじかい はりは、 さいごに すぎた \R{数字}{すう|じ}を よむ。
\end{taisetsu}

\bksubsection{3}{\R{長}{なが}い はりで 「なん\R{分}{ぷん}」}

\label{bk:1-10:fun}\R{長}{なが}い はりが \R{小}{ちい}さい \R{目}{め}もりを 1つ すすむと、 \mbox{1\R{分}{ぷん}}です。 12の ところが \mbox{0\R{分}{ぷん}}。 12から いくつ すすんだかを かぞえます。

\R{数字}{すう|じ}の あいだは 5\R{分}{ふん}。 だから、 \R{数字}{すう|じ}の ところで 5、 10、 15… と 5ずつ かぞえられます（まわりの \R{青}{あお}い \R{数}{かず}）。

\noindent\begin{minipage}[c]{0.4\linewidth}\centering
\begin{tikzpicture}[x=1mm, y=1mm]
  \bkclock{(0,0)}{17}{7}{23}
  \bkqifives{(0,0)}{17}{60}
\end{tikzpicture}
\end{minipage}\hfill
\begin{minipage}[c]{0.58\linewidth}
\begin{enumerate}[label={\bkcircled[bkTheme]{\arabic*}}, leftmargin=2\zw, itemsep=0pt, topsep=0pt]
\item \R{長}{なが}い はりの \R{前}{まえ}の \R{数字}{すう|じ}は 4。 5、 10、 15、 20で \mbox{20\R{分}{ぷん}}。
\item 4から あと \R{小}{ちい}さい \R{目}{め}もり 3つ。 21、 22、 23。
\item みじかい はりは 7と 8の あいだ。
\end{enumerate}
\par\vspace{2pt}\hspace*{2\zw}{\sffamily\bfseries\Large\color{bkTheme}\mbox{7\R{時}{じ}23\R{分}{ぷん}}}
\end{minipage}
\par\vspace{3pt}
\noindent 12から 1つずつ 23まで かぞえても、 おなじ \mbox{23\R{分}{ぷん}}に なります。 5ずつ かぞえる ほうが はやいね。\\
\R{長}{なが}い はりが さす \R{数字}{すう|じ}と 「なん\R{分}{ぷん}」の ひょうです。

\begin{center}
\begin{tblr}{colspec={Q[l,m]*{12}{Q[c,m]}}, hlines={0.6pt,bkLine}, vlines={0.6pt,bkLine}, colsep=2.2pt, rowsep=2pt,
  column{1}={bg=bkBlue!10, font=\sffamily\bfseries\footnotesize}, cells={font=\sffamily\small}}
\R{長}{なが}い はりの \R{数字}{すう|じ} & 1 & 2 & 3 & 4 & 5 & 6 & 7 & 8 & 9 & 10 & 11 & 12 \\
なん\R{分}{ぷん} & 5 & 10 & 15 & 20 & 25 & 30 & 35 & 40 & 45 & 50 & 55 & 0 \\
\end{tblr}
\end{center}
\noindent 12の ところは、 1まわりして 60\R{分}{ぷん}。 60\R{分}{ぷん}は 1\R{時間}{じ|かん}なので、 つぎの \R{時}{じ}の 0\R{分}{ぷん}に なります（⑤で くわしく やります）。

\noindent\begin{minipage}[c]{0.66\linewidth}
\R{長}{なが}い はりが 6を さす ときは \mbox{30\R{分}{ぷん}}。 \R{長}{なが}い はりは 1まわりの ちょうど \R{半分}{はん|ぶん}を すすみました（\R{右}{みぎ}の \R{色}{いろ}の ところ）。 だから \mbox{8\R{時}{じ}30\R{分}{ぷん}}を 「8\R{時}{じ}\R{半}{はん}」とも いいます。
\end{minipage}\hfill
\begin{minipage}[c]{0.32\linewidth}\centering
\begin{tikzpicture}[x=1mm, y=1mm]
  \fill[bkYellow!35] (0,0) -- (90:15) arc (90:-90:15) -- cycle;
  \bkclock{(0,0)}{15}{8}{30}
  \fill[bkYellow!35, opacity=0.55] (0,0) -- (90:14.2) arc (90:-90:14.2) -- cycle;
\end{tikzpicture}
\end{minipage}



\clearpage
\noindent{\sffamily\bfseries\color{bkTheme}よみかたの まとめ}\par
\noindent\begin{minipage}[c]{0.34\linewidth}\centering
\begin{tikzpicture}[x=1mm, y=1mm]
  \bkclock{(0,0)}{16}{6}{38}
\end{tikzpicture}
\end{minipage}\hfill
\begin{minipage}[c]{0.64\linewidth}
\begin{enumerate}[label={\bkcircled[bkTheme]{\arabic*}}, leftmargin=2\zw, itemsep=0pt, topsep=0pt]
\item みじかい はりは 6と 7の あいだ。 すぎた \R{数字}{すう|じ}は 6で、 6\R{時}{じ}。
\item \R{長}{なが}い はりの \R{前}{まえ}の \R{数字}{すう|じ}は 7で、 35\R{分}{ふん}。
\item 7から \R{小}{ちい}さい \R{目}{め}もり 3つで、 36、 37、 38。
\end{enumerate}
\par\vspace{2pt}\hspace*{2\zw}{\sffamily\bfseries\Large\color{bkTheme}\mbox{6\R{時}{じ}38\R{分}{ぷん}}}
\end{minipage}

\par\vspace{6pt}
\noindent{\sffamily\bfseries\color{bkTheme}はりの ない \R{時計}{と|けい}も ある}\par
\noindent\begin{minipage}[c]{0.6\linewidth}
でんしレンジや スマートフォンの \R{時計}{と|けい}は、 \R{数字}{すう|じ}だけで \R{時}{じ}こくを \R{見}{み}せます。 「：」の \R{左}{ひだり}が 「なん\R{時}{じ}」、 \R{右}{みぎ}が 「なん\R{分}{ぷん}」です。 \R{右}{みぎ}の 2つは、 どちらも \mbox{10\R{時}{じ}10\R{分}{ぷん}}です。
\end{minipage}\hfill
\begin{minipage}[c]{0.38\linewidth}\centering
\begin{tikzpicture}[x=1mm, y=1mm]
  \bkclock{(0,0)}{10}{10}{10}
  \fill[bkInk!85, rounded corners=1.5mm] (13,-6) rectangle (39,6);
  \node[font=\sffamily\bfseries\fontsize{17}{17}\selectfont, text=bkGreen!50] at (26,0) {10:10};
\end{tikzpicture}
\end{minipage}

\par\vspace{6pt}
\noindent{\sffamily\bfseries\color{bkTheme}ソラの 1\R{日}{にち}を \R{時計}{と|けい}で よもう}\par
\begin{bkfigure}
\begin{tikzpicture}[x=1mm, y=1mm]
  \foreach \i/\h/\m/\t/\u in {0/7/0/{おきる}/{7\R{時}{じ}},1/8/15/{いえを でる}/{8\R{時}{じ}15\R{分}{ふん}},2/12/30/{きゅうしょく}/{12\R{時}{じ}30\R{分}{ぷん}},3/3/45/{いえに かえる}/{3\R{時}{じ}45\R{分}{ふん}}} {
    \bkclock{(14+\i*33,0)}{12}{\h}{\m}
    \node[anchor=south, font=\sffamily\bfseries\footnotesize] at (14+\i*33,13.5) {\t};
    \node[anchor=north, font=\sffamily\bfseries\footnotesize, text=bkTheme] at (14+\i*33,-13.5) {\u};}
\end{tikzpicture}
\end{bkfigure}
\noindent どれも、 みじかい はりで 「なん\R{時}{じ}」、 \R{長}{なが}い はりで 「なん\R{分}{ぷん}」を よんで います。 じぶんの \R{時計}{と|けい}で たしかめて みよう。

\clearpage
\bksubsection{4}{\R{時}{じ}こくと \mbox{\R{時間}{じ|かん}}}

\label{bk:1-10:jikan}ソラは \mbox{9\R{時}{じ}}に いえを でて、 \mbox{12\R{時}{じ}}に おじいちゃんの いえに つきました。

\begin{bkfigure}
\begin{tikzpicture}[x=1mm, y=1mm]
  \fill[bkBlue!8, draw=bkBlue!50, line width=0.6pt, rounded corners=1mm] (-4,-2.5) rectangle (126,2.5);
  \foreach \h/\t in {0/8,1/9,2/10,3/11,4/12} {
    \draw[bkInk, line width=0.9pt] (8+\h*28,2.5) -- (8+\h*28,-2.5);
    \node[anchor=north, font=\sffamily\bfseries\small] at (8+\h*28,-3.5) {\t\R{時}{じ}};}
  \foreach \h in {0,...,3} \draw[bkInk!70, line width=0.5pt] (22+\h*28,2.5) -- (22+\h*28,0);
  \foreach \h in {1,2,3} {
    \draw[bkTheme, line width=1.3pt, -{Stealth[length=5pt]}] (8+\h*28+1,3.5) .. controls (8+\h*28+8,11) and (8+\h*28+20,11) .. (8+\h*28+27,3.7);
    \node[font=\sffamily\bfseries\footnotesize, text=bkTheme] at (22+\h*28,13) {\mbox{1\R{時間}{じ|かん}}};}
  \fill[bkTheme] (36,0) circle (1.6) (120,0) circle (1.6);
  \bkclock{(36,-21)}{9.5}{9}{0}
  \bkclock{(120,-21)}{9.5}{12}{0}
  \node[anchor=east, font=\sffamily\bfseries\small] at (25,-21) {でた};
  \node[anchor=east, font=\sffamily\bfseries\small] at (109,-21) {ついた};
  \node[anchor=west, font=\sffamily\bfseries\footnotesize, text=bkBlue!70!black] at (-4,7.5) {\R{時間}{じ|かん}の \R{線}{せん}};
  \node[font=\sffamily\bfseries\footnotesize, text=bkBlue!70!black, align=center] at (72,-21) {1-6の \R{数直線}{すう|ちょく|せん}と おなじ。\\ \R{右}{みぎ}へ いくほど あと。};
\end{tikzpicture}
\end{bkfigure}

\noindent \mbox{9\R{時}{じ}}や \mbox{12\R{時}{じ}}のような \mbox{\R{時計}{と|けい}}の よみを、 {\sffamily\bfseries\color{bkTheme}\R{時}{じ}こく}と いいます。\bkidx{じこく}{\R{時}{じ}こく} \R{時}{じ}こくは、 \R{時間}{じ|かん}の \R{線}{せん}の 「1つの \R{点}{てん}」です。\\
でてから つくまでの ながさ（\mbox{3\R{時間}{じ|かん}}）を、 {\sffamily\bfseries\color{bkTheme}\mbox{\R{時間}{じ|かん}}}と いいます。\bkidx{じかん}{\R{時間}{じ|かん}} \R{時間}{じ|かん}は、 2つの \R{点}{てん}の 「あいだの ながさ」です。

\noindent \mbox{\R{時間}{じ|かん}}は、 \R{時間}{じ|かん}の \R{線}{せん}で 1\R{時間}{じ|かん}ずつ かぞえても、 1-5の 「あと － まえ」の ひき\R{算}{ざん}でも もとめられます。
\bkbigshiki{12 － 9 ＝ 3}{おわりの \R{時}{じ}こく － はじめの \R{時}{じ}こく}
\noindent\bkhatenaref の A の \mbox{3\R{時}{じ}}は、 \mbox{\R{時計}{と|けい}}が さす \R{時}{じ}こく。 B の \mbox{3\R{時間}{じ|かん}}は、 でてから つくまでの ながさ （\mbox{\R{時間}{じ|かん}}） です。

\clearpage
\bksubsection{5}{\mbox{60\R{分}{ぷん}}で \mbox{1\R{時間}{じ|かん}}}

\label{bk:1-10:rokujuu}やってみようで、 \R{長}{なが}い はりが 1まわり すると、 みじかい はりは つぎの \R{数字}{すう|じ}まで すすみました。 \R{長}{なが}い はりの 1まわりは、 \R{小}{ちい}さい \R{目}{め}もり 60こ、 つまり \mbox{60\R{分}{ぷん}}。 みじかい はりの 1つ\R{分}{ぶん}は \mbox{1\R{時間}{じ|かん}}。 だから \mbox{60\R{分}{ぷん}} ＝ \mbox{1\R{時間}{じ|かん}}です。\bkidx{ふん}{\R{分}{ふん}}

\noindent 1-7では、 10 あつまると くり\R{上}{あ}がりました。 \mbox{\R{時計}{と|けい}}では、 \mbox{60\R{分}{ぷん}} あつまると \mbox{1\R{時間}{じ|かん}}に くり\R{上}{あ}がります。

\begin{bkfigure}
\begin{tikzpicture}[x=1mm, y=1mm]
  \fill[bkPaper, rounded corners=2mm] (-3,-9) rectangle (40,17);
  \foreach \i in {0,...,9} {\bkcube{(\i*3.7,11)}{3}{bkRoute3}}
  \draw[bk arrow, bkGray, line width=1pt] (17,9.5) -- (17,4);
  \foreach \i in {0,...,9} {\filldraw[fill=bkRoute3, draw=bkRoute3!60!black, line width=0.5pt] (1+\i*3.2,-1) rectangle ++(3.2,3);}
  \draw[bkRoute3!50!black, line width=1pt] (1,-1) rectangle (33,2);
  \node[anchor=north, font=\sffamily\bfseries\footnotesize] at (17,-2.5) {10こで 1つの まとまり};
  \fill[bkPaper, rounded corners=2mm] (46,-9) rectangle (132,17);
  \foreach \i in {0,...,59} {\pgfmathtruncatemacro\bkqib{mod(\i,5)==0}
    \ifnum\bkqib=1 \draw[bkInk, line width=0.7pt] (49+\i*1.33,10) -- (49+\i*1.33,14.5);
    \else\draw[bkInk!70, line width=0.35pt] (49+\i*1.33,10) -- (49+\i*1.33,12.5);\fi}
  \draw[bkInk, line width=0.7pt] (49+60*1.33,10) -- (49+60*1.33,14.5);
  \draw[bkInk, line width=0.6pt] (49,10) -- (49+60*1.33,10);
  \node[anchor=west, font=\sffamily\bfseries\footnotesize] at (49,6) {\mbox{1\R{分}{ぷん}}が 60こ};
  \draw[bk arrow, bkGray, line width=1pt] (89,8.5) -- (89,3);
  \bktape{(49,-6)}{79.8}{bkTheme}{\mbox{60\R{分}{ぷん}}で \mbox{1\R{時間}{じ|かん}}}
\end{tikzpicture}
\end{bkfigure}

\noindent{\sffamily\bfseries\color{bkTheme}80\R{分}{ぷん}は なん\R{時間}{じ|かん} なん\R{分}{ぷん}？}\par
80\R{分}{ぷん}から、 \mbox{1\R{時間}{じ|かん}}に なる \mbox{60\R{分}{ぷん}}を とりだします。 のこりは 80 － 60 ＝ 20で \mbox{20\R{分}{ぷん}}。 だから \mbox{1\R{時間}{じ|かん}20\R{分}{ぷん}}です。

\begin{bkfigure}
\begin{tikzpicture}[x=1mm, y=1mm]
  \filldraw[fill=bkTheme!25, draw=bkTheme!80!black, line width=0.8pt] (0,0) rectangle (66,7);
  \filldraw[fill=bkBlue!25, draw=bkBlue!80!black, line width=0.8pt] (66,0) rectangle (88,7);
  \node[font=\sffamily\bfseries\footnotesize] at (33,3.5) {\mbox{60\R{分}{ぷん}}};
  \node[font=\sffamily\bfseries\footnotesize] at (77,3.5) {\mbox{20\R{分}{ぷん}}};
  \bkbrace{(0,8)}{(88,8)}{\mbox{80\R{分}{ぷん}}}{1}
  \bkbrace{(0,-1)}{(66,-1)}{\mbox{1\R{時間}{じ|かん}}}{-1}
\end{tikzpicture}
\end{bkfigure}

\noindent{\sffamily\bfseries\color{bkTheme}なぜ 60に したの？}\quad 60は、 \R{半分}{はん|ぶん}の 30、 その \R{半分}{はん|ぶん}の 15、 3つに \R{分}{わ}けた 20の ように、 きりの よい \R{数}{かず}に \R{分}{わ}けやすい \R{数}{かず}です。 だから、 むかしの \R{人}{ひと}は 60を つかいました。

\clearpage
\bksubsection{6}{もっと みじかい \mbox{\R{時間}{じ|かん}}：\R{秒}{びょう}}

\label{bk:1-10:byou}\noindent\begin{minipage}[c]{0.72\linewidth}
かけっこは \mbox{1\R{分}{ぷん}}も かかりません。 こんな みじかい \mbox{\R{時間}{じ|かん}}は、 \R{秒}{びょう}で はかります。\bkidx{びょう}{\R{秒}{びょう}}
「いーち」と ゆっくり いうと、 だいたい \mbox{1\R{秒}{びょう}}。
ストップウォッチの はりは、 \R{小}{ちい}さい \R{目}{め}もり 1つで \mbox{1\R{秒}{びょう}}、 1まわりで \mbox{60\R{秒}{びょう}}です。\bkidx{すとっぷうぉっち}{ストップウォッチ} \mbox{60\R{秒}{びょう}}で \mbox{1\R{分}{ぷん}}に くり\R{上}{あ}がります。
\end{minipage}\hfill
\begin{minipage}[c]{0.26\linewidth}\centering
\begin{tikzpicture}[x=1mm, y=1mm]
  \bkqiwatch{(0,0)}{14}{12}
  \node[anchor=north, font=\sffamily\bfseries\footnotesize] at (0,-16.5) {\mbox{12\R{秒}{びょう}}};
\end{tikzpicture}
\end{minipage}

\begin{bkfigure}
\begin{tikzpicture}[x=1mm, y=1mm]
  \foreach \t/\x/\c in {{\mbox{1\R{時間}{じ|かん}}}/0/bkTheme,{\mbox{1\R{分}{ぷん}}}/48/bkBlue,{\mbox{1\R{秒}{びょう}}}/96/bkGreen} {
    \node[draw=\c, line width=1pt, fill=\c!10, rounded corners=2mm, minimum width=26mm, minimum height=9mm, font=\sffamily\bfseries\small] at (\x+13,0) {\t};}
  \foreach \x in {26,74} {\draw[bk arrow, bkGray, line width=1pt] (\x+2,0) -- (\x+20,0);
    \node[font=\sffamily\bfseries\footnotesize, text=bkGray] at (\x+11,4.5) {60に \R{分}{わ}ける};}
\end{tikzpicture}
\end{bkfigure}

\noindent{\sffamily\bfseries\color{bkTheme}どれくらいの ながさ？}\par
\noindent\hspace*{1\zw}\begin{tblr}{colspec={Q[l,m,wd=34mm]Q[l,m,wd=88mm]}, hlines={0.6pt,bkLine}, rowsep=2pt, colsep=4pt, column{1}={font=\sffamily\bfseries\small}}
だいたい \mbox{1\R{秒}{びょう}} & \R{手}{て}を 1かい たたく。 「いーち」と ゆっくり いう。 \\
だいたい \mbox{1\R{分}{ぷん}} & 60まで ゆっくり かぞえる。 \\
\mbox{45\R{分}{ふん}} & \R{小学校}{しょう|がっ|こう}の じゅぎょう 1つ\R{分}{ぶん}。 \\
\mbox{1\R{時間}{じ|かん}} & \R{長}{なが}い はりが 1まわり。 \\
\end{tblr}

\begin{taisetsu}
みじかい はりは すぎた \R{数字}{すう|じ}、 \R{長}{なが}い はりは 12から いくつ すすんだか。\\
「\R{時}{じ}こく」は \mbox{\R{時計}{と|けい}}の よみ。 「\mbox{\R{時間}{じ|かん}}」は、 はじめから おわりまでの ながさ。\\
\mbox{1\R{時間}{じ|かん}} ＝ \mbox{60\R{分}{ぷん}}\hspace{2\zw}\mbox{1\R{分}{ぷん}} ＝ \mbox{60\R{秒}{びょう}}
\end{taisetsu}

\clearpage
\bksubsection{7}{\mbox{\R{午前}{ご|ぜん}}と \mbox{\R{午後}{ご|ご}}}

\label{bk:1-10:gozen}あさ おきる \mbox{7\R{時}{じ}}と、 よる ごはんを たべる \mbox{7\R{時}{じ}}。 \mbox{\R{時計}{と|けい}}の はりは、 どちらも おなじ ところを さします。

\begin{bkfigure}
\begin{tikzpicture}[x=1mm, y=1mm]
  \bkclock{(20,0)}{13}{7}{0}
  \bkclock{(100,0)}{13}{7}{0}
  \fill[bkYellow] (20,21) circle (3.2);
  \foreach \a in {0,45,...,315} \draw[bkYellow, line width=1pt] ($(20,21)+(\a:4.4)$) -- ($(20,21)+(\a:6)$);
  \fill[bkYellow] (100,21) circle (3.5); \fill[white] (101.8,22.3) circle (3.1);
  \node[anchor=west, font=\sffamily\bfseries\small] at (36,3) {あさの \mbox{7\R{時}{じ}}};
  \node[anchor=east, font=\sffamily\bfseries\small] at (84,3) {よるの \mbox{7\R{時}{じ}}};
  \node[font=\sffamily\bfseries\footnotesize, text=bkGray] at (60,-6) {はりは おなじ};
\end{tikzpicture}
\end{bkfigure}
\noindent みじかい はりは、 \mbox{1\R{日}{にち}}に 2まわり します。 1まわりで \mbox{12\R{時間}{じ|かん}}なので、 \mbox{1\R{日}{にち}}は 12 ＋ 12 ＝ 24で \mbox{24\R{時間}{じ|かん}}です。 だから \mbox{7\R{時}{じ}}は \mbox{1\R{日}{にち}}に 2かい あります。

\begin{bkfigure}
\begin{tikzpicture}[x=1mm, y=1mm]
  \fill[bkBlue!18] (0,0) rectangle (60,7);
  \fill[bkOrange!25] (60,0) rectangle (120,7);
  \draw[bkInk!70, line width=0.6pt] (0,0) rectangle (120,7);
  \foreach \h in {0,...,24} {\pgfmathtruncatemacro\m{mod(\h,6)}\ifnum\m=0 \draw[bkInk, line width=0.9pt] (\h*5,0) -- (\h*5,-2.4);\else\draw[bkInk!60, line width=0.4pt] (\h*5,0) -- (\h*5,-1.4);\fi}
  \foreach \h/\t in {0/{0\R{時}{じ}},6/6,12/{12\R{時}{じ}},18/6,24/{12\R{時}{じ}}} {\node[anchor=north, font=\sffamily\bfseries\scriptsize] at (\h*5,-2.8) {\t};}
  \node[font=\sffamily\bfseries\small, text=bkBlue!80!black] at (30,3.5) {\R{午前}{ご|ぜん}（よなか → ひる）};
  \node[font=\sffamily\bfseries\small, text=bkOrange!85!black] at (90,3.5) {\R{午後}{ご|ご}（ひる → よなか）};
  \draw[bkTheme, line width=1pt, -{Stealth[length=3.5pt]}] (60,15) -- (60,8);
  \node[anchor=south, font=\sffamily\bfseries\footnotesize, text=bkTheme] at (60,15) {ひるの 12\R{時}{じ}（\R{正午}{しょう|ご}）};
  \draw[bkGreen!70!black, line width=1pt, -{Stealth[length=3.5pt]}] (35,-11) -- (35,-3);
  \node[anchor=north, font=\sffamily\bfseries\scriptsize, text=bkGreen!60!black] at (35,-11) {あさの 7\R{時}{じ}};
  \draw[bkGreen!70!black, line width=1pt, -{Stealth[length=3.5pt]}] (95,-11) -- (95,-3);
  \node[anchor=north, font=\sffamily\bfseries\scriptsize, text=bkGreen!60!black] at (95,-11) {よるの 7\R{時}{じ}};
\end{tikzpicture}
\end{bkfigure}
\noindent よなかの 12\R{時}{じ}から ひるの 12\R{時}{じ}までを {\sffamily\bfseries\color{bkTheme}\R{午前}{ご|ぜん}}、 ひるの 12\R{時}{じ}から よなかの 12\R{時}{じ}までを {\sffamily\bfseries\color{bkTheme}\R{午後}{ご|ご}}と いいます。\bkidx{ごぜん}{\R{午前}{ご|ぜん}}\bkidx{ごご}{\R{午後}{ご|ご}} ひるの 12\R{時}{じ}を \R{正午}{しょう|ご}とも いいます。\\
あさは \mbox{\R{午前}{ご|ぜん}7\R{時}{じ}}、 よるは \mbox{\R{午後}{ご|ご}7\R{時}{じ}}と いえば、 まちがえずに つたわります。

\clearpage
%% ---------------- 例題 ----------------
\begin{reidai}[\mbox{\R{時計}{と|けい}}を よむ]
\noindent\begin{minipage}[c]{0.7\linewidth}
この \mbox{\R{時計}{と|けい}}は、 なん\R{時}{じ} なん\R{分}{ぷん}？

\kangaekata みじかい はりは 4を すぎて、 5の \R{手前}{て|まえ}。 だから 4\R{時}{じ}。 \R{長}{なが}い はりは 7。 5、 10、 15、 20、 25、 30、 35で \mbox{35\R{分}{ふん}}。

\kotae \mbox{4\R{時}{じ}35\R{分}{ふん}}
\end{minipage}\hfill
\begin{minipage}[c]{0.28\linewidth}\centering
\begin{tikzpicture}[x=1mm, y=1mm]
  \bkclock{(0,0)}{15}{4}{35}
\end{tikzpicture}
\end{minipage}
\end{reidai}

\begin{jibun}
\noindent\begin{minipage}[c]{0.7\linewidth}
この \mbox{\R{時計}{と|けい}}は、 なん\R{時}{じ} なん\R{分}{ぷん}？\par\vspace{4pt}
\noindent\bkqijifun
\end{minipage}\hfill
\begin{minipage}[c]{0.28\linewidth}\centering
\begin{tikzpicture}[x=1mm, y=1mm]
  \bkclock{(0,0)}{14}{10}{15}
\end{tikzpicture}
\end{minipage}
\bkjans{\mbox{10\R{時}{じ}15\R{分}{ふん}}}
\end{jibun}

\begin{reidai}[1\R{分}{ぷん}まで よむ]
\noindent\begin{minipage}[c]{0.7\linewidth}
この \mbox{\R{時計}{と|けい}}は、 なん\R{時}{じ} なん\R{分}{ぷん}？

\kangaekata みじかい はりは 2と 3の あいだで 2\R{時}{じ}。 \R{長}{なが}い はりの \R{前}{まえ}の \R{数字}{すう|じ}は 9で \mbox{45\R{分}{ふん}}。 そこから \R{小}{ちい}さい \R{目}{め}もり 2つで、 46、 47。

\kotae \mbox{2\R{時}{じ}47\R{分}{ふん}}
\end{minipage}\hfill
\begin{minipage}[c]{0.28\linewidth}\centering
\begin{tikzpicture}[x=1mm, y=1mm]
  \bkclock{(0,0)}{15}{2}{47}
\end{tikzpicture}
\end{minipage}
\end{reidai}

\begin{jibun}
\noindent\begin{minipage}[c]{0.7\linewidth}
この \mbox{\R{時計}{と|けい}}は、 なん\R{時}{じ} なん\R{分}{ぷん}？\par\vspace{4pt}
\noindent\bkqijifun
\end{minipage}\hfill
\begin{minipage}[c]{0.28\linewidth}\centering
\begin{tikzpicture}[x=1mm, y=1mm]
  \bkclock{(0,0)}{14}{9}{12}
\end{tikzpicture}
\end{minipage}
\bkjans{\mbox{9\R{時}{じ}12\R{分}{ふん}}}
\end{jibun}

\begin{reidai}[かかった \mbox{\R{時間}{じ|かん}}]
\mbox{3\R{時}{じ}10\R{分}{ぷん}}に しゅくだいを はじめて、 \mbox{3\R{時}{じ}40\R{分}{ぷん}}に おわりました。 なん\R{分}{ぷん} かかった？

\noindent\begin{minipage}[c]{0.52\linewidth}
\kangaekata どちらも 3\R{時}{じ}なので、 \R{分}{ふん}だけ くらべる。 おわり － はじめ。

\shiki 40 － 10 ＝ 30\par
\kotae \mbox{30\R{分}{ぷん}}
\end{minipage}\hfill
\begin{minipage}[c]{0.46\linewidth}\centering
\begin{tikzpicture}[x=1mm, y=1mm]
  \bkclock{(0,0)}{11.5}{3}{10}
  \bkclock{(33,0)}{11.5}{3}{40}
  \node[anchor=north, font=\sffamily\bfseries\footnotesize] at (0,-12.5) {はじめ};
  \node[anchor=north, font=\sffamily\bfseries\footnotesize] at (33,-12.5) {おわり};
  \draw[bk arrow, bkBlue, line width=1.1pt] (13,0) -- (20,0);
\end{tikzpicture}
\end{minipage}
\end{reidai}

\begin{jibun}
\mbox{4\R{時}{じ}5\R{分}{ふん}}から \mbox{4\R{時}{じ}50\R{分}{ぷん}}まで あそびました。 なん\R{分}{ぷん}？
\bkshikiran
\bkjans{50 － 5 ＝ 45（1-7の ひっ\R{算}{さん}）\quad \mbox{45\R{分}{ふん}}}
\end{jibun}

\begin{reidai}[\mbox{60\R{分}{ぷん}}を \mbox{1\R{時間}{じ|かん}}に]
えんそくで \mbox{80\R{分}{ぷん}} あるきました。 なん\mbox{\R{時間}{じ|かん}} なん\R{分}{ぷん}？

\kangaekata \mbox{60\R{分}{ぷん}}で \mbox{1\R{時間}{じ|かん}}。 80から 60を とると、 80 － 60 ＝ 20。 のこりは \mbox{20\R{分}{ぷん}}。

\kotae \mbox{1\R{時間}{じ|かん}20\R{分}{ぷん}}
\end{reidai}

\begin{jibun}
\mbox{90\R{秒}{びょう}}は、 なん\R{分}{ぷん} なん\R{秒}{びょう}？\par
\noindent{\small\color{bkOrange!58!black}ヒント：\mbox{60\R{秒}{びょう}}で \mbox{1\R{分}{ぷん}}}
\bkqitwo{\R{分}{ふん}}{\R{秒}{びょう}}
\bkjans{\mbox{1\R{分}{ぷん}30\R{秒}{びょう}}（90 － 60 ＝ 30）}
\end{jibun}

\begin{reidai}[ちょうどの \R{時}{じ}こくを またぐ \mbox{\R{時間}{じ|かん}}]
\mbox{8\R{時}{じ}40\R{分}{ぷん}}から \mbox{9\R{時}{じ}10\R{分}{ぷん}}まで \R{本}{ほん}を よみました。 なん\R{分}{ぷん}？
\begin{bkfigure}
\begin{tikzpicture}[x=1mm, y=1mm]
  \fill[bkBlue!8, draw=bkBlue!50, line width=0.6pt, rounded corners=1mm] (-4,-2.5) rectangle (106,2.5);
  \foreach \i in {0,...,6} {\draw[bkInk, line width=0.6pt] (\i*17,2.5) -- (\i*17,-2.5);}
  \draw[bkInk, line width=1.4pt] (51,2.5) -- (51,-2.5);
  \foreach \i/\t in {1/{\mbox{8\R{時}{じ}40\R{分}{ぷん}}},3/{\mbox{9\R{時}{じ}}},4/{\mbox{9\R{時}{じ}10\R{分}{ぷん}}}} {
    \node[anchor=north, font=\sffamily\bfseries\footnotesize] at (\i*17,-3.5) {\t};}
  \fill[bkBlue] (17,0) circle (1.6) (68,0) circle (1.6);
  \draw[bkBlue, line width=1.3pt, -{Stealth[length=5pt]}] (18,3.5) .. controls (28,10) and (40,10) .. (50.5,3.7);
  \node[font=\sffamily\bfseries\footnotesize, text=bkBlue] at (34,11) {\mbox{20\R{分}{ぷん}}};
  \draw[bkTheme, line width=1.3pt, -{Stealth[length=5pt]}] (52,3.5) .. controls (56,9) and (63,9) .. (67.5,3.7);
  \node[font=\sffamily\bfseries\footnotesize, text=bkTheme] at (60,10.5) {\mbox{10\R{分}{ぷん}}};
\end{tikzpicture}
\end{bkfigure}
\kangaekata \R{時}{じ}が ちがうので、 ちょうどの \mbox{9\R{時}{じ}}で 2つに \R{分}{わ}ける。 9\R{時}{じ}までは 60 － 40 ＝ 20で \mbox{20\R{分}{ぷん}}。 9\R{時}{じ}から 9\R{時}{じ}10\R{分}{ぷん}までは \mbox{10\R{分}{ぷん}}。

\shiki 20 ＋ 10 ＝ 30 \kotaenext \mbox{30\R{分}{ぷん}}
\end{reidai}

\begin{jibun}
\mbox{1\R{時}{じ}50\R{分}{ぷん}}から \mbox{2\R{時}{じ}25\R{分}{ふん}}まで そうじを しました。 なん\R{分}{ぷん}？\par
\noindent{\small\color{bkOrange!58!black}ヒント：まず \mbox{2\R{時}{じ}}まで なん\R{分}{ぷん}？}
\bkshikiran
\bkjans{60 － 50 ＝ 10、 10 ＋ 25 ＝ 35\quad \mbox{35\R{分}{ふん}}}
\end{jibun}

\begin{reidai}[\R{時}{じ}こくを もとめる]
\mbox{7\R{時}{じ}40\R{分}{ぷん}}の \mbox{30\R{分}{ぷん}} あとは、 なん\R{時}{じ} なん\R{分}{ぷん}？
\begin{bkfigure}
\begin{tikzpicture}[x=1mm, y=1mm]
  \fill[bkBlue!8, draw=bkBlue!50, line width=0.6pt, rounded corners=1mm] (-4,-2) rectangle (90,2);
  \foreach \i in {0,...,5} {\draw[bkInk, line width=0.6pt] (\i*17,2) -- (\i*17,-2);}
  \draw[bkInk, line width=1.4pt] (51,2) -- (51,-2);
  \foreach \i/\t in {1/{\mbox{7\R{時}{じ}40\R{分}{ぷん}}},3/{\mbox{8\R{時}{じ}}},4/{\mbox{8\R{時}{じ}10\R{分}{ぷん}}}} {
    \node[anchor=north, font=\sffamily\bfseries\scriptsize] at (\i*17,-2.8) {\t};}
  \fill[bkBlue] (17,0) circle (1.4);
  \draw[bkBlue, line width=1.2pt, -{Stealth[length=4pt]}] (18,2.8) .. controls (28,8) and (40,8) .. (50.5,3);
  \node[font=\sffamily\bfseries\scriptsize, text=bkBlue] at (34,9) {\mbox{20\R{分}{ぷん}}};
  \draw[bkTheme, line width=1.2pt, -{Stealth[length=4pt]}] (52,2.8) .. controls (56,7.5) and (63,7.5) .. (67.5,3);
  \node[font=\sffamily\bfseries\scriptsize, text=bkTheme] at (60,9) {\mbox{10\R{分}{ぷん}}};
\end{tikzpicture}
\end{bkfigure}
\kangaekata 40 ＋ 30 ＝ 70だと \mbox{7\R{時}{じ}70\R{分}{ぷん}}に なって しまう。 ちょうどの \mbox{8\R{時}{じ}}で \R{分}{わ}ける（れいだい5と おなじ）。\\
\mbox{8\R{時}{じ}}まで あと 60 － 40 ＝ 20で \mbox{20\R{分}{ぷん}}。 30 － 20 ＝ 10で、 のこり \mbox{10\R{分}{ぷん}}。

\kotae \mbox{8\R{時}{じ}10\R{分}{ぷん}}
\end{reidai}

\begin{jibun}
\mbox{3\R{時}{じ}45\R{分}{ふん}}の \mbox{30\R{分}{ぷん}} あとは、 なん\R{時}{じ} なん\R{分}{ぷん}？\par
\noindent{\small\color{bkOrange!58!black}ヒント：あと \mbox{15\R{分}{ふん}}で \mbox{4\R{時}{じ}}}
\bkqijifun
\bkjans{\mbox{4\R{時}{じ}15\R{分}{ふん}}（60 － 45 ＝ 15、 30 － 15 ＝ 15）}
\end{jibun}

\clearpage
%% ---------------- p.104 まちがえやすいところ ----------------
\bksection{tsumazuki}{まちがえやすい ところ}

\begin{tsumazuki}{\mbox{8\R{時}{じ}}？ \mbox{9\R{時}{じ}}？}
\noindent\begin{minipage}[c]{0.78\linewidth}
\bkbatsu\enspace \mbox{9\R{時}{じ}50\R{分}{ぷん}}\hspace{3\zw}\bkmaru\enspace \mbox{8\R{時}{じ}50\R{分}{ぷん}}\par
まだ 9に ついて いないので、 \mbox{8\R{時}{じ}}です。
\end{minipage}\hfill
\begin{minipage}[c]{0.2\linewidth}\centering
\begin{tikzpicture}[x=1mm, y=1mm]
  \bkclock{(0,0)}{11.5}{8}{50}
\end{tikzpicture}
\end{minipage}
\end{tsumazuki}

\begin{tsumazuki}{\mbox{4\R{分}{ぷん}}？ \mbox{20\R{分}{ぷん}}？}
\noindent\begin{minipage}[c]{0.78\linewidth}
\bkbatsu\enspace \mbox{2\R{時}{じ}4\R{分}{ぷん}}\hspace{3\zw}\bkmaru\enspace \mbox{2\R{時}{じ}20\R{分}{ぷん}}\par
4は \mbox{20\R{分}{ぷん}}。 5、 10、 15、 20と かぞえよう。
\end{minipage}\hfill
\begin{minipage}[c]{0.2\linewidth}\centering
\begin{tikzpicture}[x=1mm, y=1mm]
  \bkclock{(0,0)}{11.5}{2}{20}
\end{tikzpicture}
\end{minipage}
\end{tsumazuki}

\begin{tsumazuki}{60で くり\R{上}{あ}げる（1-7）}
\mbox{4\R{時}{じ}50\R{分}{ぷん}}の \mbox{20\R{分}{ぷん}} あとは？\par
\bkbatsu\enspace \mbox{4\R{時}{じ}70\R{分}{ぷん}}\hspace{3\zw}\bkmaru\enspace \mbox{5\R{時}{じ}10\R{分}{ぷん}}\par
\mbox{70\R{分}{ぷん}}は ありません。 \mbox{60\R{分}{ぷん}}で \mbox{1\R{時間}{じ|かん}}に くり\R{上}{あ}げます。
\end{tsumazuki}

\begin{tsumazuki}{「\R{時}{じ}」と 「\R{時間}{じ|かん}」を まぜる}
\mbox{9\R{時}{じ}}から \mbox{11\R{時}{じ}}まで あそんだ。\hspace{1\zw}\bkbatsu\enspace \mbox{2\R{時}{じ}}\hspace{2\zw}\bkmaru\enspace \mbox{2\R{時間}{じ|かん}}\par
あそんだ ながさは 「\R{時間}{じ|かん}」で いいます。
\end{tsumazuki}



\clearpage
%% ---------------- p.105 つかわれているところ ----------------
\begin{butsuri}
\R{時}{じ}こくと \mbox{\R{時間}{じ|かん}}は、 のりものや きかいを うごかして います。
\begin{bkfigure}
\begin{tikzpicture}[x=1mm, y=1mm]
  % 電車の時こく表
  \fill[white, rounded corners=3mm] (0,0) rectangle (39,38);
  \fill[bkInk!85, rounded corners=1mm] (4,21) rectangle (35,35);
  \node[anchor=west, text=bkYellow, font=\sffamily\bfseries\fontsize{8}{9}\selectfont] at (5.5,31.2) {でる\enspace\mbox{8\R{時}{じ}15\R{分}{ふん}}};
  \node[anchor=west, text=bkGreen!40, font=\sffamily\bfseries\fontsize{8}{9}\selectfont] at (5.5,24.8) {つく\enspace\mbox{8\R{時}{じ}45\R{分}{ふん}}};
  \fill[bkRoute4, rounded corners=1.5mm] (5,10) rectangle (34,18);
  \foreach \i in {0,1,2} {\fill[white, rounded corners=0.6mm] (8+\i*8.5,13) rectangle ++(6,3.6);}
  \fill[bkInk!80] (11,10) circle (1.5) (28,10) circle (1.5);
  \draw[bkInk, line width=0.8pt] (3,8.4) -- (36,8.4);
  \node[font=\sffamily\bfseries\footnotesize] at (19.5,4) {\R{電車}{でん|しゃ}};
  % しんごう
  \begin{scope}[shift={(42,0)}]
  \fill[white, rounded corners=3mm] (0,0) rectangle (39,38);
  \fill[bkInk!80, rounded corners=1.5mm] (8,10) rectangle (21,35);
  \fill[bkRed!30] (14.5,28.5) circle (4.4);
  \fill[bkGreen] (14.5,16.3) circle (4.4);
  \bkperson{(14.5,12.8)}{7}{white}
  \foreach \k in {0,...,4} {\fill[\ifnum\k<3 bkGreen\else bkGray!40\fi, rounded corners=0.4mm] (25,11+\k*4.6) rectangle ++(6,3.4);}
  \node[font=\sffamily\bfseries\footnotesize] at (19.5,4) {しんごう};
  \end{scope}
  % ロボットそうじき
  \begin{scope}[shift={(84,0)}]
  \fill[white, rounded corners=3mm] (0,0) rectangle (39,38);
  \fill[bkGray!25] (3,10) rectangle (36,12);
  \filldraw[fill=bkNavy!80, draw=bkInk!70, line width=0.6pt] (17.5,19) ellipse (11 and 6);
  \fill[bkNavy!55] (17.5,20.5) ellipse (8 and 3.6);
  \fill[bkGreen] (17.5,21.5) circle (1.2);
  \foreach \r in {4,7,10} {\draw[bkOrange, line width=0.9pt] ($(28.5,19)+(-35:\r)$) arc (-35:35:\r);}
  \node[font=\sffamily\bfseries\footnotesize] at (19.5,4) {ロボットそうじき};
  \node[font=\sffamily\bfseries\fontsize{8}{9}\selectfont, text=bkOrange!80!black] at (19.5,32) {\mbox{1\R{秒}{びょう}}に なんかいも};
  \end{scope}
\end{tikzpicture}
\end{bkfigure}
\begin{itemize}[label={\color{bkPurple}●}, leftmargin=1.4\zw, itemsep=0pt, topsep=0pt]
\item \mbox{\R{電車}{でん|しゃ}}は、 \R{時}{じ}こく\R{表}{ひょう}の \R{時}{じ}こくに でて、 きまった \mbox{\R{時間}{じ|かん}}で つぎの えきに つきます。
\item しんごうは、 あおの \mbox{\R{時間}{じ|かん}}を \R{秒}{びょう}で きめて、 いろを かえます。
\item ロボットそうじきは、 \mbox{1\R{秒}{びょう}}に なんかいも まわりを しらべて、 まえと くらべて どれだけ かわったかを みて います。
\end{itemize}
\end{butsuri}

\bkfukidashi{delta}{\mbox{\R{時間}{じ|かん}}は、 ものが どれだけ かわったかを はかる\\「ものさし」。 「\mbox{\R{変化}{へん|か}}の \mbox{\R{路線}{ろ|せん}}」は、 ここも とおるよ！}

\begin{bkfigure}
\begin{tikzpicture}[x=1mm, y=1mm]
  \draw[bkRoute1, line width=3pt] (10,0) -- (128,0);
  \node[anchor=west, inner sep=0pt] at (0,12) {\bkrouteicon[1.4em]{1}\hspace{0.3em}{\sffamily\bfseries\small\color{bkRoute1!80!black}\R{変化}{へん|か}の \R{路線}{ろ|せん}}};
  \foreach \x/\c/\t/\s in {10/bkVol1/{1-5}/{ひき\R{算}{ざん}},39.5/bkVol1/{1-10}/{\R{時間}{じ|かん}},69/bkVol1/{1-14}/{1つ\R{分}{ぶん}},98.5/bkVol1/{1-18}/{おれ\R{線}{せん}グラフ},128/bkVol2/{2-3}/{{\ltjsetparameter{jacharrange={+200}}＋－×÷}}} {
    \fill[white, draw=\c, line width=1.6pt] (\x,0) circle (2.6);
    \node[anchor=base, font=\sffamily\bfseries\footnotesize, text=\c!85!black] at (\x,-8.4) {\t};
    \node[anchor=base, font=\sffamily\footnotesize, text=bkInk] at (\x,-15.4) {\s};}
  \node[fill=bkTheme, text=white, rounded corners=1mm, font=\sffamily\bfseries\fontsize{8}{9}\selectfont, inner sep=1.5pt] at (39.5,5.5) {いま ここ};
\end{tikzpicture}
\end{bkfigure}


\clearpage
%% ---------------- p.106〜107 もんだい ----------------
\bkendproblems

\bkproblemgroup{きほん}{1}

\bkproblem \mbox{\R{時計}{と|けい}}を よみましょう。
\begin{bkfigure}
\begin{tikzpicture}[x=1mm, y=1mm]
  \foreach \i/\h/\m in {0/7/0,1/2/30,2/10/43,3/8/50} {
    \bkclock{(17+\i*33,0)}{11}{\h}{\m}
    \node[font=\sffamily\bfseries\small] at (2+\i*33,11) {（\the\numexpr\i+1\relax）};}
\end{tikzpicture}
\par\vspace{2pt}
\noindent\begin{tblr}{colspec={*{4}{Q[c,m,wd=33mm]}}, colsep=0pt, rowsep=0pt}
\bkkakikomi[8mm]\,\R{時}{じ} & \bkkakikomi[8mm]\,\R{時}{じ}\,\bkkakikomi[8mm]\,\R{分}{ふん} & \bkkakikomi[8mm]\,\R{時}{じ}\,\bkkakikomi[8mm]\,\R{分}{ふん} & \bkkakikomi[8mm]\,\R{時}{じ}\,\bkkakikomi[8mm]\,\R{分}{ふん} \\
\end{tblr}
\end{bkfigure}
\bkpans{\mbox{（1）7\R{時}{じ}}\quad \mbox{（2）2\R{時}{じ}30\R{分}{ぷん}}\quad \mbox{（3）10\R{時}{じ}43\R{分}{ぷん}}\quad \mbox{（4）8\R{時}{じ}50\R{分}{ぷん}}}

\bkproblem えんそくで、 \mbox{9\R{時}{じ}}に がっこうを でて、 \mbox{11\R{時}{じ}}に こうえんに つきました。 なん\mbox{\R{時間}{じ|かん}} かかった？
\bkshikiran
\bkpans{11 － 9 ＝ 2\quad \mbox{2\R{時間}{じ|かん}}}

\bkproblem \mbox{6\R{時}{じ}15\R{分}{ふん}}に ごはんを たべはじめて、 \mbox{6\R{時}{じ}40\R{分}{ぷん}}に たべおわりました。 なん\R{分}{ぷん} たべて いた？
\bkshikiran
\bkpans{40 － 15 ＝ 25（1-7の ひっ\R{算}{さん}）\quad \mbox{25\R{分}{ふん}}}

\bkproblem \bkkakikomi[8mm]に かずを かきましょう。
\par\vspace{2pt}
\noindent\hspace*{1\zw}（1）\mbox{1\R{時間}{じ|かん}} ＝ \bkkakikomi\,\R{分}{ふん}\hspace{3\zw}（2）\mbox{1\R{分}{ぷん}} ＝ \bkkakikomi\,\R{秒}{びょう}\par\vspace{4pt}
\noindent\hspace*{1\zw}（3）\mbox{70\R{分}{ぷん}} ＝ \bkkakikomi\,\mbox{\R{時間}{じ|かん}}\,\bkkakikomi\,\R{分}{ふん}
\bkpans{\mbox{（1）60}\quad \mbox{（2）60}\quad \mbox{（3）1\R{時間}{じ|かん}10\R{分}{ぷん}}}

\bkproblem \mbox{2\R{時}{じ}40\R{分}{ぷん}}から \mbox{3\R{時}{じ}20\R{分}{ぷん}}まで、 こうえんで あそびました。 なん\R{分}{ぷん} あそんだ？
\bkshikiran
\bkpans{60 － 40 ＝ 20、 20 ＋ 20 ＝ 40\quad \mbox{40\R{分}{ぷん}}}

\bkproblemgroup{チャレンジ}{1}

\bkproblem \begin{minipage}[t]{0.68\linewidth}
\mbox{\R{時計}{と|けい}}を みましょう。 この \mbox{30\R{分}{ぷん}} あとは、 なん\R{時}{じ} なん\R{分}{ぷん}？\par
{\small\color{bkOrange!58!black}ヒント：\mbox{11\R{時}{じ}}まで あと なん\R{分}{ぷん}？}
\bkqijifun
\end{minipage}\hfill
\begin{minipage}[t]{0.26\linewidth}\raggedleft
\begin{tikzpicture}[x=1mm, y=1mm, baseline={([yshift=-3mm]current bounding box.north)}]
  \bkclock{(0,0)}{13}{10}{50}
\end{tikzpicture}
\end{minipage}
\bkpans{\mbox{11\R{時}{じ}20\R{分}{ぷん}}}

\bkproblem \mbox{\R{電子}{でん|し}}レンジで、 おべんとうを \mbox{100\R{秒}{びょう}} あたためます。 なん\R{分}{ぷん} なん\R{秒}{びょう}？
\bkqitwo{\R{分}{ふん}}{\R{秒}{びょう}}
\bkpans{\mbox{1\R{分}{ぷん}40\R{秒}{びょう}}}
\begin{furikaeri}
\begin{checklist}
\item \mbox{\R{時計}{と|けい}}を よんで、 \R{時}{じ}こくが いえる。\hfill{\footnotesize\sffamily\color{bkTheme}→\pageref{bk:1-10:shikumi}ページ}
\item \R{時}{じ}こくと \mbox{\R{時間}{じ|かん}}の ちがいが わかる。\hfill{\footnotesize\sffamily\color{bkTheme}→\pageref{bk:1-10:jikan}ページ}
\item 60で くり\R{上}{あ}がる けいさんが できる。\hfill{\footnotesize\sffamily\color{bkTheme}→\pageref{bk:1-10:rokujuu}ページ}
\end{checklist}
チェックが つかなかったら、 →の ページに もどろう。
\end{furikaeri}

\clearpage
